


\documentclass[12]{article}
\usepackage[utf8]{inputenc}
%\usepackage[latin1]{inputenc}%\textbf{}
\usepackage[spanish, es-tabla]{babel} %convierte comandos en ingles al espa\~{n}ol. 
\usepackage[ ]{graphicx} %paquete para poder incluir graficas
%\DeclareGraphicsExtensions{.png,.pdf,.jpg,.jpeg,.xps}
\usepackage[right=2cm,left=3cm,top=2.5cm,bottom=2.0cm,headsep=0cm,footskip=0.5cm]{geometry} %margenes del texto
%\usepackage{anysize} %Soporte para el comando de margenes
%\marginsize{3cm}{2.cm}{2cm}{2cm} % Controla losmárgenes{izquierda}{derecha}{arriba}% {abajo}. 
\setlength{\parindent}{0cm}%anula la sangrí en todo el texto

\title{\huge{\textbf{Escriba el t\'itulo de acuerdo al objetivo}}}%Comando para el título

\author{\small \textit{{K. A. Silva-Fonseca, A. M. Verdeza-Caballero, M. F, Mej\'ia-Castilla, M. P. Castro V\'elez.}}}%Introducción de los nombres de los integrantes

%CUERPO DEL DOCUMENTO

\begin{document}

\maketitle

\begin{center}
\rule[0mm]{150mm}{0.1mm}% Para dibujar una linea horizontal de
		                        %[elevación]{longitud}{grosor}
\end{center}


\begin{abstract}
\noindent %anula la sangria de este parrafo
Esta es una prueba para editar un articulo en \LaTeX en Linux Mint. Mint es un sistema operativo basado en Linux Ubuntu, el cual es m\'as f\'acil de instalar. En este SO se pueden instalar los siguientes editores de textos \LaTeX: Kile, TeXworks, TeXstudio y funcionan muy bien. Para que el editor funcione muy bien, se debe instalar el paquete Texlive que es equivalente al Miktex para windows.\end{abstract}

\begin{center}
 \rule[0mm]{150mm}{0.1mm}% Para dibujar una linea horizontal de [elevacón]{longitud}{grosor}
\end{center}



\section{Introducci\'on}

Esta es una prueba para editar un articulo en LaTeX en Linux Mint. Mint es un sistema operativo basado en Linux Ubuntu, el cual es m\'as f\'acil de instalar. En este SO se pueden  instalar los siguientes editores de textos LaTeX: Kile, TeXworks, TeXstudio y funcionan muy bien. Para que el editor funcione muy bien, se debe instalar el paquete Texlive que es equivalente al Miktex para windows.

\section{Discusi\'on te\'orica}

Estuvimos hoy en el laboratorio aprendiendo a editar en Latex
\noindent
Para escribir ecuaciones en la misma linea usa $ y = x^{2}$
\noindent
Para editar ecuaciones en otra linea y centradas sin numeraci\'on, escribes la aecuacion entre doble s\'imbolo de pesos
$$ y = x^{2}$$
\noindent
Para editar ecuaciones centradas y numeradas usas los comandos

\begin{equation}
	y = x^{2}
\end{equation}

%\newpage % esta isntruccion salta de pagina cuando sea necesario. para desactivarla escribe el simbolo % antes de la barra \

\section{M\'etodo Experimental}

Para llevar a cabo este experimento, se uso una torrre con dos alambres conectados a alta tensi\'on, entre los que cae un bal\'in...

Nota: no liste los elementes usados, describa como los utiliz\'o.

\section{Resultados y discusi\'on}

Para introducir tablas

\begin{table}[h]
	\centering
		\caption{Datos de posici\'on tiempo}
  	\begin{tabular}{|l|r|c|c|} % Barras para hacer linea vertical. l=izq, r = der, c=centrado
		\hline % hace la linea horiontal
		Tiempo (s) & Posici\'on(cm) & Log t & log x \\
		\hline
		0         & 0          & -     & -\\
		\hline
		1        & 1         & 0 & 0 \\ %doble barra invert, salta linea (copiar y pegar para obtener las filas)
		\hline
		2         & 4          & 0,301     & 0,602\\
		\hline
		3          & 9          & 0,477     & 0,954\\
		\hline
		4          & 16          & 0,602    & 1,204\\
		\hline
		5         & 25         & 0,699     & 1,398\\
		\hline
		\end{tabular}
	
	\label{tab:Datos1}
\end{table}

\newpage %Este comando lo pueden desactivar si no se necesita

\noindent
Para introducir figuras usa el paquete  \textbackslash usepackage{graphix} en el pre\'ambulo 
y el comando \textbackslash includegraphics[scale\\=x.x]{archivo. jpg}.

\begin{figure}[h]
  \begin{center}
    \includegraphics[scale=0.7]{figuraxy2.jpg}
    \caption{Gr\'afico de la ecuaci\'on parab\'olica $y = x^{2}$} 
    \label{fig:grafico}
  \end{center}
\end{figure}
Los archivos de las figuras deben guardarse en la misma carpeta del archivo .tex, de otra forma tiene que escribir la ruta de la carpeta donde guarda el archivo jpg.\\

Para obtener las constantes de la ecuaci\'on, se usa el m\'etodo de linealizaci\'on descrito en la referencia \cite{guia}.


\section{Conclusiones}

%\section*{Referencias}

\begin{thebibliography}{10}

{%apertura
\setlength{\parindent}{0cm} %Estas instrucciones anulan la sangr´\'ia en los párrafos que esten dentro de los corchetes

\bibitem{baird} D. C. Baird, Experimentaci\'on: Una introducci\'on a la teor\'ia de mediciones
y al dise\~{n}o de experimentos. Prentice-Hall Hispanoamericana, S. A. Segunda
Edici\'on. 1991

\bibitem{eecoral} E. E. Coral-Escobar. Dise\~{n}o y construcci\'on de un medidor de conductividad t\'ermica de materiales aislantes, tipo cer\'amicos. Revista colombiana de F\'isica, vol. 38, No. 3, 2006).


\bibitem{guia} E. E. Coral-Escobar. Gu\'ia para An\'alisis de Experimentos. Universidad del Atl\'antico. Barranquilla, 2011
}%cierre

\end{thebibliography}

\end{document}
